\subsection{Experimental Setup and Protocols}
\label{sec:experimental_protocols}

The experiments address three questions: whether recorded PX4 telemetry supports runtime anomaly detection and conditional fault-domain diagnosis; whether hierarchical SHAP aggregation provides traceable, prediction-relevant diagnostic evidence; and whether commit-specific uORB relations can convert topic evidence into useful software-module suspect rankings under controlled source mutations. Table~\ref{tab:protocol} summarizes the datasets, splits, evaluation roles, and permitted outputs. Detailed preprocessing and model definitions are given in Section 3 and are not repeated here; the measures used across all experiments are defined centrally in Section~\ref{sec:evaluation_metrics}.

UAV-SEAD is the only dataset used to fit the supervised Stage 1 and Stage 2 models. The fixed chronological split provides the primary reproduction protocol, the 14-window purged split tests sensitivity to temporal proximity, and leave-log-out assigns complete flights to partitions to test unseen-log generalization. ALFA is used only for zero-shot binary detection because its labels are not mapped to the four UAV-SEAD fault domains. Uncategorized UAV-SEAD logs support only flight-level weak-label screening because they have no temporal anomaly ranges. Controlled PX4 replay is separate from real-flight model fitting and provides the only source-mutation ground truth for module-level ranking.

Baseline comparisons distinguish detector performance from the structure of diagnosis. Direct Five-Class is the principal structural comparator to the complete Cascade, while Random Forest, majority, and stratified-random references retain their corresponding supervised tasks. Resource-limited CATCH and linear VAR(1) feasibility results, together with their computational budgets and evaluation scope, are documented separately in Supplementary Material Section S6.

For the supervised real-flight models, thresholds and early stopping are selected on validation data only. Replay residual calibration instead uses development healthy pairs as specified below. Five model seeds are used for the principal real-flight experiments. Seed standard deviations describe fitted-model variation, whereas 95\% cluster-bootstrap intervals describe sampling variation: complete flight logs are resampled for UAV-SEAD, sequences for ALFA, and flights for Uncategorized screening. Absolute and paired real-flight analyses use 1000 repetitions with seed 20260825. New replay results are descriptive counts over three source flights per group; old intervals computed from incomplete replay inputs are not transferred to corrected results.

Corrected replay fixes PX4 commit \texttt{\seqsplit{82aa24adfca29321cfd1209e287eab6c2b16780e}} and Ubuntu 20.04 WSL, without Gazebo, ROS, QGroundControl, or NuttX. The initial three development source logs and three additional transfer logs each have four mutation conditions (Commander, EKF2, INAV, and Land Detector), with a healthy reference, fault target, and independent normal target. Fixed mutation activation at 30 s is used for evaluation only. Source logs lack \texttt{ekf2\_timestamps}, so EKF2 uses the generic publication adapter with \texttt{ekf2 start}, without the dedicated \texttt{-r} handshake. Input completion, clock alignment, finite estimator outputs, and the original effect checks are reported separately from predictive scores.

The first transfer protocol froze its residual parameters before evaluating the three additional sources. Those sources had appeared in the historical P2 scaler data, so the isolation concerns replay calibration, not a globally untouched preprocessing pipeline. Subsequent common-support and calibration analyses use the same already examined transfer sources and are exploratory. A two-reference sensitivity adds 24 healthy replays across the same six source logs, using the new runs as normal targets and the two older healthy outputs as references. These repeats add no independent source flights. Supplementary Material Sections S1 and S5 specify source identities, data roles, and all controls.

\begin{table}[H]
\caption{Evaluation protocols and data partitions.}
\label{tab:protocol}
\centering
\footnotesize
\renewcommand{\arraystretch}{1.14}
\begin{tabularx}{\textwidth}{@{}>{\RaggedRight\arraybackslash}p{2.15cm} >{\RaggedRight\arraybackslash}p{2.85cm} >{\RaggedRight\arraybackslash}p{5.05cm} >{\RaggedRight\arraybackslash}X@{}}
\toprule
\textbf{Protocol} & \textbf{Purpose} & \textbf{Data and evaluation set} & \textbf{Isolation / scope}\\
\midrule
Fixed chronological & Primary reproduction & 1389 UAV-SEAD logs; 152,340/32,138/34,059 train/validation/test windows; five seeds & Chronological windows; overlapping source logs may span partitions\\
Purged & Temporal-proximity sensitivity & 133,353/8,536/17,966 windows; 14-window boundary purge; five seeds & Removes overlap near split boundaries\\
Leave-log-out & Unseen-flight generalization & 970/206/213 complete train/validation/test logs; 149,774/37,324/31,439 windows & No flight log crosses a partition\\
ALFA zero-shot & External binary transfer & 46 sequences; 5273 windows; 16 shared channels & Frozen UAV-SEAD detector and validation threshold; no four-class transfer\\
Uncategorized & Weak-label screening & 139 Uncategorized versus 898 pure-Normal logs & Flight-level metrics only; no temporal event labels\\
Paired PX4 replay & Controlled extension & 3 development + 3 transfer sources; 12 fault targets per group; separate healthy targets/references & Frozen transfer followed by explicitly exploratory analyses; source identities in Supplement S1\\
\bottomrule
\end{tabularx}
\par\vspace{2pt}{\scriptsize All UAV-SEAD protocols use 96-sample windows, stride 8, horizon 12, 87 channels, and 1566 derived features. Sampling intervals use 1000 flight- or sequence-cluster bootstrap repetitions. Corrected replay counts are descriptive; source-flight dependence is retained. Frozen model configurations and per-run records are described in the Data Availability Statement.}
\end{table}
